% Práctica 05 | Responsable: Miguel
\section{Clientes y contactos}

Se toma como base la conversión del informe de la Práctica~04. Los
atributos multivaluados de contacto se representan en relaciones
independientes y las partidas conservan su contexto identificador.
Todas las columnas del módulo son obligatorias y carecen de valor
predeterminado. Las tablas se crean en \texttt{tablas.sql}; la nulabilidad,
los dominios y las llaves se agregan mediante \texttt{ALTER TABLE}.

\begingroup
\small
\setlength{\tabcolsep}{5pt}
\renewcommand{\arraystretch}{1.2}

\subsection{Cliente}

\noindent\begin{tabular}{@{}p{.25\textwidth}p{.20\textwidth}p{.10\textwidth}p{.36\textwidth}@{}}
\hline
\textbf{Columna} & \textbf{Tipo} & \textbf{NULL} & \textbf{Dominio} \\
\hline
\texttt{idCliente} & \texttt{integer} & No & Entero positivo; identificador único del cliente. \\
\texttt{fechaNacimiento} & \texttt{date} & No & Fecha válida y finita, no posterior a la fecha actual al insertar o actualizar la fila. \\
\texttt{nombre} & \texttt{varchar(50)} & No & Hasta 50 caracteres; al menos uno distinto de espacio en blanco. \\
\texttt{apellidoMaterno} & \texttt{varchar(50)} & No & Hasta 50 caracteres; al menos uno distinto de espacio en blanco. \\
\texttt{apellidoPaterno} & \texttt{varchar(50)} & No & Hasta 50 caracteres; al menos uno distinto de espacio en blanco. \\
\texttt{sexo} & \texttt{varchar(50)} & No & Descriptor no vacío; sin catálogo cerrado en el modelo. \\
\hline
\end{tabular}

La llave primaria es \texttt{idCliente}. El nombre completo se descompone
en nombre y apellidos. Se permiten acentos, espacios y signos en estos
campos, pero se rechazan cadenas vacías o formadas únicamente por
espacios en blanco. La edad se obtiene de la fecha de nacimiento y la
categoría VIP se calcula con la regla del caso; no se almacenan como
columnas redundantes.

\subsection{CorreoCliente}

\noindent\begin{tabular}{@{}p{.25\textwidth}p{.20\textwidth}p{.10\textwidth}p{.36\textwidth}@{}}
\hline
\textbf{Columna} & \textbf{Tipo} & \textbf{NULL} & \textbf{Dominio} \\
\hline
\texttt{idCliente} & \texttt{integer} & No & Entero positivo que referencia a un cliente existente. \\
\texttt{correoElectrónico} & \texttt{varchar(100)} & No & Hasta 100 caracteres; una arroba, parte local y dominio no vacíos, sin espacios en blanco. \\
\hline
\end{tabular}

La llave primaria es \texttt{(idCliente, correoElectrónico)}. La FK
\texttt{CorreoCliente.idCliente} $\to$ \texttt{Cliente.idCliente}
garantiza la existencia del propietario. Un cliente puede registrar
varios correos y distintos clientes pueden compartir la misma dirección;
lo que no puede repetirse es la pareja completa. La validación del correo
es básica: no comprueba la existencia del buzón ni toda su sintaxis.

\subsection{TelefonoCliente}

\noindent\begin{tabular}{@{}p{.25\textwidth}p{.20\textwidth}p{.10\textwidth}p{.36\textwidth}@{}}
\hline
\textbf{Columna} & \textbf{Tipo} & \textbf{NULL} & \textbf{Dominio} \\
\hline
\texttt{idCliente} & \texttt{integer} & No & Entero positivo que referencia a un cliente existente. \\
\texttt{teléfono} & \texttt{char(10)} & No & Exactamente diez dígitos de 0 a 9. \\
\hline
\end{tabular}

La llave primaria es \texttt{(idCliente, teléfono)} y la FK referencia
a \texttt{Cliente.idCliente}. El teléfono se guarda como cadena para
conservar ceros iniciales. Un cliente puede registrar varios teléfonos;
la misma cadena puede pertenecer a más de un cliente sin violar la PK.

\clearpage
\subsection{Tarjeta}

\noindent\begin{tabular}{@{}p{.25\textwidth}p{.20\textwidth}p{.10\textwidth}p{.36\textwidth}@{}}
\hline
\textbf{Columna} & \textbf{Tipo} & \textbf{NULL} & \textbf{Dominio} \\
\hline
\texttt{idCliente} & \texttt{integer} & No & Entero positivo que identifica al cliente propietario. \\
\texttt{numeroTarjeta} & \texttt{integer} & No & Entero positivo conservado como atributo de la tarjeta. \\
\texttt{costoAdquisición} & \texttt{integer} & No & Importe entero no negativo en pesos. \\
\texttt{fechaCompra} & \texttt{date} & No & Fecha válida y finita del calendario. \\
\texttt{saldo} & \texttt{integer} & No & Saldo entero no negativo en pesos. \\
\texttt{estado} & \texttt{varchar(50)} & No & Descriptor de hasta 50 caracteres, no vacío; sin catálogo cerrado. \\
\texttt{puntosAcumulados} & \texttt{integer} & No & Entero no negativo. \\
\texttt{CURP} & \texttt{char(18)} & No & Identificador de un cajero existente. \\
\hline
\end{tabular}

\paragraph{Identificación y relación con Cliente.}
La llave primaria es \texttt{idCliente}, que también referencia a
\texttt{Cliente.idCliente}. De esta manera cada tarjeta pertenece a un
cliente y cada cliente tiene como máximo una tarjeta. El informe de la
Práctica~04 conserva \texttt{numeroTarjeta} como dato, sin hacerlo único
globalmente ni incluirlo en una llave redundante. Por ello no se agrega
\texttt{UNIQUE(numeroTarjeta)} ni se propaga ese atributo a Partida.

\paragraph{Vendedor y dominios.}
La FK \texttt{Tarjeta.CURP} $\to$ \texttt{Cajero.CURP} exige que
el vendedor exista en el subtipo Cajero; no basta con que exista en
Empleado. Se conserva el dominio de su CURP por esa referencia.
Los importes y puntos son enteros no negativos: cero es un valor válido.
No se asignan saldos, costos, puntos o estados predeterminados que el
informe no establece. La fecha de compra debe ser finita; el informe no
fija una restricción respecto de la fecha actual para esta columna.

\paragraph{Participación y operaciones.}
La FK desde Tarjeta y su PK no obligan a que todos los clientes tengan
tarjeta. La participación total de Cliente en Poseer requiere validación
adicional entre tablas. Asimismo, los \texttt{CHECK} impiden almacenar
valores negativos, pero no sustituyen la comprobación y actualización
conjunta de saldos y puntos al realizar recargas, partidas y canjes.

\clearpage
\subsection{Partida}

\noindent\begin{tabular}{@{}p{.25\textwidth}p{.20\textwidth}p{.10\textwidth}p{.36\textwidth}@{}}
\hline
\textbf{Columna} & \textbf{Tipo} & \textbf{NULL} & \textbf{Dominio} \\
\hline
\texttt{idCliente} & \texttt{integer} & No & Entero positivo que referencia a la tarjeta utilizada. \\
\texttt{idMaquina} & \texttt{integer} & No & Entero positivo que referencia a una máquina existente. \\
\texttt{idPartida} & \texttt{integer} & No & Entero positivo; discriminante dentro del contexto tarjeta y máquina. \\
\texttt{fecha} & \texttt{date} & No & Fecha válida y finita del calendario. \\
\texttt{hora} & \texttt{timetz} & No & Hora válida del día con desplazamiento de zona horaria. \\
\texttt{puntosGanados} & \texttt{integer} & No & Entero no negativo. \\
\texttt{puntajeObtenido} & \texttt{integer} & No & Entero no negativo. \\
\hline
\end{tabular}

La llave primaria compuesta es
\texttt{(idCliente, idMaquina, idPartida)}, en ese orden. Identifica a
la entidad débil dentro de sus propietarios. Un mismo valor de
\texttt{idPartida} puede reutilizarse si cambia la tarjeta o la máquina;
solo se rechaza la repetición de los tres componentes simultáneamente.

Se establecen las FK \texttt{Partida.idCliente} $\to$
\texttt{Tarjeta.idCliente} y \texttt{Partida.idMaquina} $\to$
\texttt{Máquina.idMaquina}. La primera exige una tarjeta existente, no
únicamente un cliente. El propietario se obtiene a través de Tarjeta.
Las referencias obligatorias garantizan que toda partida tenga ambos
participantes; una tarjeta o máquina puede existir sin partidas.

La hora se conserva como \texttt{timetz}, de acuerdo con la unificación
de tipos del informe anterior. En las inserciones de prueba se proporciona
el desplazamiento de zona explícitamente. La bonificación VIP es un
atributo derivado y no se añade como columna. Los límites de puntos o
puntaje que dependan de un juego concreto deben validarse al registrar
la operación; aquí se conserva únicamente su dominio no negativo.

\clearpage
\subsection{CorreoEmpleado}

\noindent\begin{tabular}{@{}p{.25\textwidth}p{.20\textwidth}p{.10\textwidth}p{.36\textwidth}@{}}
\hline
\textbf{Columna} & \textbf{Tipo} & \textbf{NULL} & \textbf{Dominio} \\
\hline
\texttt{CURP} & \texttt{char(18)} & No & Identificador de un empleado existente. \\
\texttt{correoElectrónico} & \texttt{varchar(50)} & No & Hasta 50 caracteres; una arroba, partes no vacías y sin espacios en blanco. \\
\hline
\end{tabular}

La llave primaria es \texttt{(CURP, correoElectrónico)}. La FK
\texttt{CorreoEmpleado.CURP} $\to$ \texttt{Empleado.CURP}
conserva la existencia y el dominio del propietario. Un empleado puede
tener varios correos y una dirección puede compartirse entre empleados.
La comprobación de formato es la misma que en CorreoCliente; se conserva
la longitud de 50 caracteres indicada para este atributo en el informe.

\subsection{TelefonoEmpleado}

\noindent\begin{tabular}{@{}p{.25\textwidth}p{.20\textwidth}p{.10\textwidth}p{.36\textwidth}@{}}
\hline
\textbf{Columna} & \textbf{Tipo} & \textbf{NULL} & \textbf{Dominio} \\
\hline
\texttt{CURP} & \texttt{char(18)} & No & Identificador de un empleado existente. \\
\texttt{teléfono} & \texttt{char(10)} & No & Exactamente diez dígitos; se conservan ceros iniciales. \\
\hline
\end{tabular}

La llave primaria es \texttt{(CURP, teléfono)} y la FK referencia a
\texttt{Empleado.CURP}. Se permite más de un teléfono por empleado y
compartirlo entre propietarios distintos. Los contactos se relacionan con
Empleado y quedan disponibles para todos sus subtipos, sin duplicarlos
en Gerente, Cajero, EncargadoDePremios o Técnico.

\subsection{Integridad referencial e integración}

Las ocho FK utilizan \texttt{ON UPDATE NO ACTION} y
\texttt{ON DELETE NO ACTION}. Se rechazan cambios que dejen referencias
sin destino. Esta política es una decisión de implementación, pues el
informe no fijaba acciones de borrado. Las restricciones no se declaran
diferibles y se comprueban al terminar cada sentencia.

El contrato del modelo sigue pendiente. Esta implementación conserva
los nombres y las llaves de la conversión del informe de la Práctica~04,
que difiere del diagrama original en Tarjeta y sus operaciones. Los
identificadores se escriben entre comillas dobles para conservar
mayúsculas y acentos. Para integrar, el módulo~02 debe confirmar
\texttt{"Máquina"} con \texttt{"idMaquina"} de tipo \texttt{integer};
el módulo~05 debe confirmar \texttt{"Empleado"} y \texttt{"Cajero"},
ambas con \texttt{"CURP"} de tipo \texttt{char(18)} y clave primaria
o única.

\subsection{Comprobación de las restricciones}

\texttt{pruebas.sql} se ejecuta con \texttt{psql} en una base de pruebas.
Crea un esquema aislado y dependencias mínimas para Empleado, Cajero y
Máquina; carga los cuatro archivos de implementación y utiliza datos
sintéticos. Comprueba las 29 columnas obligatorias, los dominios, las
siete PK y las ocho FK. También comprueba contactos compartidos, varios
contactos por propietario y discriminantes de partida reutilizados en
contextos distintos. Cada error esperado se aísla y verifica mediante
su \texttt{SQLSTATE} y, cuando corresponde, el nombre de la restricción.
Un resultado distinto detiene la ejecución. Se muestran los resultados
y se revierte toda la transacción.

Esta validación del módulo no sustituye la ejecución del DDL completo
con los cinco módulos reales ni las evidencias de restauración del respaldo.
\endgroup
